\documentclass[11pt,a4paper]{article}
\usepackage[margin=1in]{geometry}
\usepackage{amsmath,booktabs,hyperref,float}

\title{Iteration 019: Stochastic Weight Averaging}
\author{Autoresearch Loop}
\date{April 8, 2026}

\begin{document}
\maketitle

\begin{abstract}
Stochastic Weight Averaging (SWA) with combined loss matches the best result
at $r=0.378$, tying iter017. SWA finds a flatter minimum but does not
improve peak performance, suggesting the loss landscape is already smooth
in this regime.
\end{abstract}

\section{Hypothesis}
SWA (Izmailov et al., 2018) averages weights from multiple points along the
training trajectory, finding solutions in flatter regions of the loss landscape
that generalize better. This could improve cross-subject transfer.

\section{Method}
Train with the iter017 recipe for 100 epochs, then collect SWA averages over
epochs 100--150 with a cyclical learning rate. Final model uses the SWA-averaged
weights with updated batch normalization statistics.

\section{Results}
\begin{table}[H]
\centering
\begin{tabular}{lrrr}
\toprule
Model & Mean $r$ & Std $r$ & SNR (dB) \\
\midrule
Combined loss (iter017) & 0.378 & 0.078 & 0.61 \\
SWA + combined loss & 0.378 & 0.078 & 0.61 \\
\bottomrule
\end{tabular}
\end{table}

\section{Analysis}
SWA matches but does not exceed iter017, suggesting the FIR model with CF init
already converges to a broad minimum. The loss landscape for this linear-ish
problem is relatively convex, so weight averaging provides no additional benefit.
The 0.378 plateau may represent a fundamental limit of linear spatio-temporal
filtering on this narrowband data.

Breaking past 0.378 likely requires either: (1) input-adaptive models that handle
cross-subject differences (attention, MoE), (2) broader frequency content, or
(3) fundamentally different data augmentation strategies.

\section{Next}
Iterations 020--023: InstanceNorm, Mixture of Experts, and transformer architectures
to explore input-adaptive approaches.

\end{document}
